
\SourceLeaf{476}

\Needspace{5\baselineskip}\begin{center}CONTENTS.\end{center}
\BookRule{0.15\linewidth}

\ContentsHeading{\textsc{Dedication}}{p. 3.}

\ContentsHeading{\textsc{Preface}}{p. 5.}

\ContentsHeading{\textsc{Section I.} Congruence of numbers in general}{p. 9.}

\ContentsEntry{Congruent numbers, moduli, residues and nonresidues, Article 1 f.}

\ContentsEntry{Least residues, 4.}

\ContentsEntry{Elementary propositions on congruent numbers, 5.}

\ContentsEntry{Some applications, 12.}

\ContentsHeading{\textsc{Section II.} Congruences of the first degree}{p. 14.}

\ContentsEntry{Preliminary theorems on prime numbers, factors, etc., Article 13.}

\ContentsEntry{Solution of congruences of the first degree, 26.}

\ContentsEntry{Finding a number congruent to given residues for given moduli, 32.}

\ContentsEntry{Linear congruences involving several unknowns, 37.}

\ContentsEntry{Various theorems, 38.}

\ContentsHeading{\textsc{Section III.} Residues of powers}{p. 38.}

\ContentsEntry{The residues of the terms of a geometric progression beginning with unity form a periodic sequence, Article 45.}

\ContentsEntry{\textit{Prime moduli are considered first.}}

\ContentsEntry{Setting the modulus \(=p\), the number of terms in the period divides \(p-1\), Article 49.}

\ContentsEntry{Fermat's theorem, 50.}

\ContentsEntry{How many numbers correspond to periods whose number of terms is a given divisor of \(p-1\), Article 52.}

\ContentsEntry{Primitive roots, bases, indices, 57.}

\ContentsEntry{The algorithm of indices, 58.}

\ContentsEntry{The roots of the congruence \(x^n\equiv A\), Article 60.}

\ContentsEntry{The relation between indices in different systems, 69.}

\ContentsEntry{Bases suited to particular uses, 72.}

\ContentsEntry{A method of finding primitive roots, 73.}

\ContentsEntry{Various theorems on periods and primitive roots, 75.}

\ContentsEntry{(Wilson's theorem, 76).}

\SourceLeaf{477}

\ContentsEntry{\textit{Moduli that are powers of prime numbers}, Article 82.}

\ContentsEntry{\textit{Moduli that are powers of two}, 90.}

\ContentsEntry{\textit{Moduli composed of several primes}, 92.}

\ContentsHeading{\textsc{Section IV.} Congruences of the second degree}{p. 73.}

\ContentsEntry{Quadratic residues and nonresidues, Article 94.}

\ContentsEntry{Whenever the modulus is prime, the number of residues smaller than it equals the number of nonresidues, 96.}

\ContentsEntry{Whether a composite number is a residue or a nonresidue of a given prime depends on the nature of its factors, 98.}

\ContentsEntry{Composite moduli, 100.}

\ContentsEntry{A general criterion for whether a given number is a residue or a nonresidue of a given prime, 106.}

\ContentsEntry{\textit{Investigations of primes for which given numbers are residues or nonresidues}, 107.}

\ContentsEntry{The residue \(-1\), Article 108.}

\ContentsEntry{The residues \(+2\) and \(-2\), Article 112.}

\ContentsEntry{The residues \(+3\) and \(-3\), Article 117.}

\ContentsEntry{The residues \(+5\) and \(-5\), Article 121.}

\ContentsEntry{On \(\pm7\), Article 124.}

\ContentsEntry{Preparation for the general investigation, 125.}

\ContentsEntry{A general (\textit{fundamental}) theorem is established by induction, and conclusions are deduced from it, 130.}

\ContentsEntry{A rigorous proof of this theorem, 135.}

\ContentsEntry{An analogous method of proving the theorem of Article 114, 145.}

\ContentsEntry{Solution of the general problem, 146.}

\ContentsEntry{Linear forms containing all primes for which any given number is either a residue or a nonresidue, 147.}

\ContentsEntry{The work of others on these investigations, 151.}

\ContentsEntry{Congruences of the second degree that are not pure, 152.}

\ContentsHeading{\textsc{Section V.} Forms and indeterminate equations of the second degree}{p. 120.}

\ContentsEntry{The aim of the investigation; definition and notation of forms, Article 153.}

\ContentsEntry{Representation of numbers; the determinant, 154.}

\ContentsEntry{Values of the expression \(\surd(bb-ac)\) (mod. \(M\)) to which the representation of the number \(M\) by the form \((a,b,c)\) belongs, 155.}

\ContentsEntry{A form containing another, or contained in another; proper and improper transformations, 157.}

\ContentsEntry{Proper and improper equivalence, 158.}

\ContentsEntry{Opposite forms, 159, adjacent forms, 160.}

\ContentsEntry{Common divisors of the coefficients of forms, 161.}

\ContentsEntry{The relation among all similar transformations of a given form into a given form, 162.}

\ContentsEntry{Ambiguous forms, 163.}

\ContentsEntry{A theorem on the case where a form is contained in another both properly and improperly, 164.}

\ContentsEntry{General results on representations of numbers by forms and their relation to transformations, 166.}

\ContentsEntry{\textit{Forms of negative determinant}, 171.}

\ContentsEntry{Special applications to decomposing numbers into two squares, into a square and twice a square, and into a square and three times a square, 182.}

\ContentsEntry{\textit{Forms of positive nonsquare determinant}, 183.}

\SourceLeaf{478}

\ContentsEntry{\textit{Forms of square determinant}, Article 206.}

\ContentsEntry{Forms contained in others to which they are nevertheless not equivalent, 213.}

\ContentsEntry{\textit{Forms of determinant \(0\)}, Article 215.}

\ContentsEntry{The general solution in integers of all indeterminate equations of the second degree involving two unknowns, 216.}

\ContentsEntry{Historical notes, 222.}

\Needspace{5\baselineskip}\begin{center}\textit{Further investigations of forms.}\end{center}

\ContentsEntry{Distribution of forms of a given determinant into classes, Article 223; of classes into orders, 226.}

\ContentsEntry{Division of orders into genera, 228.}

\ContentsEntry{\textit{Composition of forms}, 234.}

\ContentsEntry{Composition of orders, 245, genera, 246, classes, 249.}

\ContentsEntry{For a given determinant, each genus of the same order contains equally many classes, 252.}

\ContentsEntry{Comparison of the numbers of classes contained in the individual genera of different orders, 253.}

\ContentsEntry{The number of ambiguous classes, 257.}

\ContentsEntry{At least half of all characters that can be assigned for a given determinant cannot have corresponding properly primitive genera (positive for negative determinants), 261.}

\ContentsEntry{A second proof of the fundamental theorem and the other theorems concerning the residues \(-1,\allowbreak +2,\allowbreak -2\), 262.}

\ContentsEntry{A closer determination of the half of the characters to which genera cannot correspond, 263.}

\ContentsEntry{A special method of decomposing primes into two squares, 265.}

\Needspace{5\baselineskip}\begin{center}\textit{A digression containing a treatment of ternary forms}, Article 266 ff.\end{center}

\ContentsEntry{\textit{Some applications to the theory of binary forms.}}

\ContentsEntry{Finding a form whose duplication produces a given binary form of the principal genus, 286.}

\ContentsEntry{Genera do in fact correspond to all characters except those found impossible in Articles 262, 263, 287 III.}

\ContentsEntry{The theory of decomposing both numbers and binary forms into three squares, 288.}

\ContentsEntry{Proof of Fermat's theorems that every integer can be decomposed into three triangular numbers or four squares, 293.}

\ContentsEntry{Solution of the equation \(axx+byy+czz=0\), Article 294.}

\ContentsEntry{The method by which the distinguished Le Gendre treated the fundamental theorem, 296.}

\ContentsEntry{Representation of zero by arbitrary ternary forms, 299.}

\ContentsEntry{The general solution in rational quantities of indeterminate equations of the second degree involving two unknowns, 300.}

\ContentsEntry{The average number of genera. 301; of classes, 302.}

\ContentsEntry{A special algorithm of properly primitive classes; regular and irregular determinants, etc., 305.}

\ContentsHeading{\textsc{Section VI.} Various applications of the preceding investigations}{p. 380.}

\ContentsEntry{Decomposition of fractions into simpler ones, Article 309.}

\ContentsEntry{Conversion of common fractions into decimals, 312.}

\ContentsEntry{Solution of the congruence \(xx\equiv A\) by the method of exclusion, 319.}

\ContentsEntry{Solution of the indeterminate equation \(mxx+nyy=A\) by exclusions, 323.}

\SourceLeaf{479}

\ContentsEntry{Another method of solving the congruence \(xx\equiv A\) for the case where \(A\) is negative, Article 327.}

\ContentsEntry{Two methods of distinguishing composite numbers from primes and finding the factors of the former, 329.}

\ContentsHeading{\textsc{Section VII.} Equations defining divisions of the circle}{p. 412.}

\ContentsEntry{The investigation is reduced to the simplest case, where the number of parts into which the circle is to be divided is prime, Article 336.}

\ContentsEntry{Equations for trigonometric functions of arcs that are one part or several parts of the whole circumference; reduction of trigonometric functions to roots of the equation \(x^n-1=0\), Article 337.}

\ContentsEntry{\textit{Theory of the roots of this equation} (where \(n\) is assumed to be prime).}

\ContentsEntry{Omitting the root \(1\), the remaining roots \((\Omega)\) are contained in the equation \(X=x^{n-1}+x^{n-2}+\mathrm{etc.}+x+1=0\).}

\ContentsEntry{The function \(X\) cannot be resolved into factors of lower degree in which all coefficients are rational, 341.}

\ContentsEntry{The aim of the following investigations is stated, 342.}

\ContentsEntry{All roots \(\Omega\) are distributed into certain classes (periods), 343.}

\ContentsEntry{Various theorems on these periods, 344 ff.}

\ContentsEntry{The solution of the equation \(X=0\) is built on these investigations, Article 352.}

\ContentsEntry{Examples for \(n=19\), where the task is reduced to two cubic equations and one quadratic, and for \(n=17\), where it is reduced to four quadratics, Articles 353. 354.}

\Needspace{4\baselineskip}\ContentsEntry{\textit{Further investigations of this subject.}}

\ContentsEntry{Sums with an even number of terms are real quantities, 355.}

\ContentsEntry{The equation defining the distribution of the roots \(\Omega\) into two periods, 356.}

\ContentsEntry{Proof of the theorem mentioned in Section IV, 357.}

\ContentsEntry{The equation for distributing the roots \(\Omega\) into three periods, 358.}

\ContentsEntry{Reduction to pure equations of the equations by which the roots \(\Omega\) are found, 359.}

\ContentsEntry{\textit{Application of the preceding investigations to trigonometric functions.}}

\ContentsEntry{A method of identifying the angles to which the individual roots \(\Omega\) correspond, 361.}

\ContentsEntry{Tangents, cotangents, secants and cosecants are derived from sines and cosines without division, 362.}

\ContentsEntry{A method of successively lowering the degrees of equations for trigonometric functions, 363.}

\ContentsEntry{Divisions of the circle that can be carried out by quadratic equations or by geometric constructions, 365.}

\ContentsHeading{\textsc{Additions}}{p. 465.}

\ContentsHeading{\textsc{Tables}}{p. 467.}

\BookRule{0.3\linewidth}

\Needspace{5\baselineskip}\begin{center}\textsc{Printed at Goslar by E. W. G. Kircher.}\end{center}
\clearpage
